5. Discussion and policy implications



\subsection{What the international evidence says, and how our results fit}



The weak association between labor market conditions and homicides documented in Section 4 is the modal finding of the verified literature, not an anomaly. In developed countries, unemployment effects are robust for property crime and small or null for homicide: meta-analytic evidence across 63 aggregate U.S. studies finds stronger unemployment effects on property than on violent crime (Chiricos, 1987); U.S. state panels imply an unemployment elasticity of roughly 0.15 for property crime and 0.04 for violent crime (Raphael \& Winter-Ebmer, 2001); Swedish data show unemployment raising thefts but not homicides (Edmark, 2005); low wages raise property crime without effects on violence (Machin \& Meghir, 2004); youth unemployment raises robberies in France (Fougère et al., 2009); and the wages of unskilled young men matter more than the unemployment rate in U.S. counties (Gould et al., 2002). Reviews conclude that labor incentives operate mainly on property crime (Draca \& Machin, 2015), and that labor market improvement did not drive the U.S. crime decline of the 1990s (Levitt, 2004).


For developing countries the evidence is more mixed but equally instructive. Cross-country panels including Latin America find that inequality and unemployment raise homicides and robbery, and that growth reduces them (Fajnzylber et al., 2002a), with income inequality -- rather than poverty -- the robust determinant of violent crime and strong persistence of crime levels (Fajnzylber et al., 2002b). Quasi-experimental work finds that Brazilian regions more exposed to trade liberalization experienced simultaneous increases in unemployment, homicides, and property crime (Dix-Carneiro et al., 2018), and that job loss raises crime, including violent crime, with unemployment insurance mitigating it substantially (Britto et al., 2022). Development reduces crime through higher opportunity costs and stronger state deterrence capacity (Soares, 2004). Yet the region's homicide spikes are consistently attributed to illegal markets and armed-group dynamics rather than aggregate unemployment: kingpin decapitation fragmented cartels and raised violence in Mexico (Dell, 2015; Calderón et al., 2015); cocaine supply shocks transmit violence across countries (Castillo et al., 2020); drug violence spreads spatially toward weakly institutionalized municipalities (Osorio, 2015); coffee price collapses -- which reduce licit income -- intensify conflict in Colombia (Dube \& Vargas, 2013); coca cultivation finances armed groups without raising household income (Angrist \& Kugler, 2008); and criminal groups use violence strategically against the state (Lessing, 2015) while providing order and protection that can keep violence low until state crackdowns break it (Lessing, 2021). Gangs' territorial control directly distorts local labor markets: residents of gang-dominated territories in El Salvador earn roughly USD 350 less per month and skilled workers migrate (Melnikov et al., 2025); childhood exposure to illegal labor markets raises the probability of adult criminal careers (Sviatschi, 2022a); and maras evolved from street gangs to transnational protection rackets (Cruz, 2010).


Our results fit the literature precisely where it is most informative. Ecuador's labor profile -- low unemployment, high underemployment, low adequate employment -- is the profile in which labor effects operate through the quality of licit alternatives rather than the unemployment rate (Fajnzylber et al., 2002a; Britto et al., 2022). The weak cross-sectional correlation (r $\leq$ 0.34) and the insignificant coefficients are what the literature predicts for homicides driven by organized crime, and the timing evidence -- escalation from 2021, record levels before the 2024 decrees, flat event study -- points to the channels emphasized regionally: illicit markets, armed-group fragmentation, and state capacity. Reverse causality also runs the opposite way: crime destroys employment and investment, with a documented macroeconomic impact in Ecuador (IMF, 2024), so any simple reading of labor--homicide correlations is doubly ambiguous.



\subsection{Policy implications with explicit evidence labels}



The following recommendations translate the findings into actions for senior officials (Presidency, Ministry of the Interior, Ministry of Labor, National Planning Secretariat), with explicit evidence labels; none claims a causal effect for any policy that has not been evaluated, and institutionalizing that standard is itself the core recommendation. No impact evaluation for Ecuador measuring the effect of labor or security programs on homicides exists in the indexed literature.


\textit{(i) Data infrastructure: make causal evaluation possible. [Evidence: high.]} The binding constraint on every causal question in this paper is the absence of a quarterly provincial labor panel: only one ENEMDU cross-section (2024-II) is available, which makes it impossible to separate labor effects on crime from the reverse channel (IMF, 2024) or to build credible counterfactuals. The highest-return investments are: publishing ENEMDU with provincial representativeness and quarterly periodicity (chained series from 2018); consolidating interoperable administrative records (police, prosecution, sentencing, prisons, IESS affiliation, transfer beneficiaries) with stable definitions, since definition changes and under-registration undermine causal attribution (Escaño et al., 2025); publishing homicide statistics at canton and weekly disaggregation; and creating a technical evidence unit for security with a mandate to evaluate exceptional-regime interventions and publish results regardless of sign. With 24 provinces, few-cluster inference methods are required (Bertrand et al., 2004; MacKinnon \& Webb, 2017), and pre-tests of parallel trends do not validate designs (Roth, 2022).


\textit{(ii) Security policy: prioritize focused policing and state presence in high-lethality territories; avoid generalized mano dura. [Evidence: high for focalization; high against generalized mano dura.]} Focused police presence has real but modest and local deterrent effects: doubled patrols in Minneapolis hot spots reduced crime calls by 6--13\% (Sherman \& Weisburd, 1995); fixed police posts in Buenos Aires cut car theft by roughly 75\% on protected blocks, with no appreciable spillover a block or two away (Di Tella \& Schargrodsky, 2004); meta-analytic evidence confirms small but significant hot-spots effects (Braga et al., 2014); and police elasticity estimates in U.S. cities are negative (Chalfin \& McCrary, 2018; Evans \& Owens, 2007; Klick \& Tabarrok, 2005). The most direct warning for Ecuador is the largest Latin American policing RCT: doubling patrols on 1,919 streets in Bogotá produced no significant crime reductions and displacement to neighboring streets, with confidence intervals ruling out total reductions above 2\% (Blattman et al., 2021). Generalized mano dura, by contrast, was ineffective or counterproductive in Central America (Jütersonke et al., 2009; Wolf, 2012): national time-series evidence for El Salvador (2003--2022) shows mano dura phases with approximately 54\% more monthly homicides than truce phases, with about 17,767 homicides attributable to the repressive period (Escaño et al., 2025); unconditional truce negotiations made violence a bargaining asset (Jones \& Lloyd, 2024); and leadership decapitation without territorial consolidation fragments groups and escalates violence (Dell, 2015; Calderón et al., 2015; Sviatschi, 2022b). In favelas with consolidated criminal governance, police presence without legitimacy produces lethal violence (Magaloni et al., 2020), and the long-run objective is rebuilding the state's monopoly of violence with clear rules and accountability (Acemoglu et al., 2013). Concrete actions: define the country's top lethality hot spots with problem-oriented patrol and quarterly targets for robbery, extortion, and homicide; accompany the exceptional response with a territorial-consolidation plan and civilian oversight; monitor displacement toward neighboring provinces (Osorio, 2015; Blattman et al., 2021); and pair policing with municipal services. Our event study neither supports nor refutes the decrees' effect; no claim about the exceptional regime can be made from these data.


\textit{(iii) Labor and prevention policy: expand youth employment programs and conditional cash transfers in high-lethality cantons, with evaluation built in and with labor-insertion targets -- not homicide targets. [Evidence: high for labor effects; low-to-medium for effects on homicides.]} The best-evidenced regional youth training programs -- Jóvenes en Acción in Colombia (Attanasio et al., 2011) and Juventud y Empleo in the Dominican Republic (Card et al., 2011; Ibarrarán et al., 2019) -- show real but moderate labor effects (+19.6\% income and +6.8 percentage points of formal salaried employment among women in Jóvenes en Acción), and none measures homicides; the Liberia experiment with high-risk men reduced hours in illicit activities without measuring homicides (Blattman \& Annan, 2016). There is no indexed evidence that youth employment reduces homicides in organized-crime settings; promising it would overstate what science can support. What does have support: conditional cash transfers are the social policy with the strongest causal evidence on crime in the region -- Bolsa Família coverage reduced urban crime in São Paulo through income and peer effects (Chioda et al., 2016) -- making the Bono de Desarrollo Humano the natural candidate for an evaluation replicating that design; licit-income shocks raise conflict, recruitment, and homicides in developing countries (Dube \& Vargas, 2013; Dix-Carneiro et al., 2018); and protecting youth in gang-controlled territories -- keeping them in school and in licit work -- is the most direct route to preventing recruitment (Melnikov et al., 2025; Sviatschi, 2022a; Cruz, 2010). Concrete actions: scale training-plus-insertion programs for 18--29-year-olds in the highest-lethality cantons with a capital component and accompaniment (Blattman \& Annan, 2016), with randomized or quasi-experimental evaluation from design; have the Planning Secretariat evaluate the BDH's effect on crime in high-coverage cantons following Chioda et al. (2016); design school-based recruitment-prevention programs without stigmatizing neighborhoods (Wolf, 2012); and report formal-insertion and school-retention targets, never homicide reductions attributed to these programs.


\textit{(iv) Institutionalize impact evaluation of security interventions. [Evidence: high.]} The exceptional-regime decrees should be evaluated with the methods validated in this project -- event study with province-and-period fixed effects and wild cluster bootstrap inference (Bertrand et al., 2004; MacKinnon \& Webb, 2017), Goodman-Bacon decomposition and Sun--Abraham estimates as diagnostics (Goodman-Bacon, 2021; Sun \& Abraham, 2021), synthetic control with placebo inference (Abadie et al., 2010), and cautious reading of pre-tests (Roth, 2022) -- with mandatory publication of results regardless of sign. Management-by-results with inter-institutional coordination, as documented for Pernambuco's Pact for Life (Ratton et al., 2014), and the regional prevention menus synthesized by the IDB and the World Bank (Jaitman, 2017; Chioda, 2017) offer the best-supported governance models. Institutional rule: every new security or employment intervention should present an ex-ante evaluation design (expected outcome, comparator, required data) together with its decree or regulation, and the legislature should require this for approving exceptional regimes and security budgets.



\subsection{What cannot be concluded}



Three statements must remain explicit. First, this paper finds no evidence either way on the exceptional regime: a null event study with a single national treatment date cannot estimate the decrees' effect. Second, the labor market results are associations; they neither prove that employment programs would fail to reduce violence nor support claims that they would. Third, the synthetic control for Esmeraldas is a descriptive case-study datum (minimum placebo p $\approx$ 1/24), not an evaluation of any intervention. Policy discourse should follow a golden rule: present findings with their limits, cite evaluations when they exist, and distinguish what is known, what is suspected, and what is unknown.